% Introduction -- written from scratch on 2026-10-02 after the 2026-09-07 rejection.
% Rewritten 2026-10-04 after the funnel of thesis/style_guide_aida.md (her ch. 1): one step per
% paragraph, claims only. Paragraph 1 is fixed. Evidence cut from here and still owed to the
% State of the Art: 196 patients, left main bifurcation angle 35.5 to 178 degrees
% \cite{temov2016bifurcation}; 136 autopsy hearts, no visible narrowing until plaque fills about 40\%
% of the cross-section \cite{glagov1987remodeling}. The 127-patient curvature result is already in
% clinical_background/05_hemodynamics.tex.
% 2026-10-04: limitations paragraph trimmed to a problem statement; the evidence (Han, Tello Ayala,
% Glagov, Temov) now lives in literature_review.tex. "considers only" softened to "rests on", since
% chan2024orfan shows imaging already refines risk in referred patients. shen2026geometry replaced
% by kashyap2022tortuosity and zebicmihic2024index (Shen excluded thesis-wide, user 2026-10-04).
% 2026-10-05: paragraph 4 reordered to match the review (measurement = Sec. 2, timing = Sec. 3);
% zebicmihic2024index swapped for cui2017bifurcation so every cited study is discussed in the review.
% No forward references in this chapter (user).

\chapter{Introduction}
\label{ch:introduction}
Cardiovascular disease is the leading cause of death worldwide, accounting for an estimated 19.2
million deaths in 2023 \cite{stark2025gbd}.
Coronary artery disease (CAD) is the largest single contributor, responsible for nearly half of
these deaths \cite{stark2025gbd}.
It develops when plaque builds up in the walls of the coronary arteries, which supply the heart
muscle with blood.
The plaque slowly narrows the arteries and can, in extreme cases, rupture, forming a clot that
blocks blood flow and causes a heart attack \cite{libby2019atherosclerosis}.

Plaque does not form evenly along the coronary tree, even though risk factors such as high blood
pressure, high cholesterol and smoking reach every segment through the blood.
Instead, it tends to collect where the drag of the blood flow on the vessel wall is low or
oscillating \cite{chatzizisis2007ess}.
This drag is known as wall shear stress, and it is low mainly on the outer wall just after a
bifurcation and on the inner wall of a bend \cite{asakura1990flow}.
Where plaque forms is therefore tied to the geometry of the tree, which, unlike shear, can be
measured directly from an image.

Coronary geometry varies considerably between individuals \cite{temov2016bifurcation}, so some
trees contain more extensive low-shear regions than others.
Predisposition to plaque may therefore depend not only on systemic risk factors but also on the
anatomy of the individual's own arteries.
Standard risk assessment, however, rests on the former, combining age, sex, smoking, blood
pressure and cholesterol into a single score \cite{score2_2021}.

Whether geometry adds to such a score remains unknown.
Studies relating geometry to plaque have defined geometry inconsistently, and the definition can
decide whether a relationship is found \cite{cui2017bifurcation,kashyap2022tortuosity}.
Most also measured geometry in referred patients whose plaque was already present, so the geometry
may be a consequence of the disease rather than its cause \cite{han2022plaque}.

The aim of this thesis is to test whether coronary geometry, measured before any symptoms of heart
disease, can categorize individuals by their risk of developing plaque.
Rather than referred patients, this work uses adults selected at random from the population of
greater Copenhagen by the Copenhagen General Population Study (CGPS) and scanned for research with
coronary computed tomography angiography (CCTA) \cite{fuchs2023cgps}.
A longitudinal subset of these individuals was scanned a second time 10 years later.
The work proceeds in two stages.
First, an automatic pipeline is developed that segments the coronary tree, assigns each vessel to
its coronary artery and computes branching angles, branch radius ratios and vessel curvature.
Second, these features are used to predict which individuals develop new plaque at a second scan.
Two risk models are compared: a logistic regression on summary features and a transformer that reads
the full curvature profile along each artery.
% TODO (author): baseline open until the CGPS variables are checked. If every SCORE2 marker is
% available the baseline is all of them; otherwise name the subset used.
Both are judged by how much they improve on a baseline model of conventional risk factors, and both are fitted
separately for each coronary artery, since shear acts locally.

Together, these stages establish an automated workflow that measures coronary geometry at population
scale and tests its link to plaque.
If geometry does carry risk, the workflow can identify individuals at increased risk from CCTA scans
that are already acquired.
If it does not, the result indicates that the localization of plaque within a tree does not translate
into a difference in risk between individuals, and the workflow remains a tool for research on
coronary anatomy and CAD.
